\begin{textsample}{POS Dim 1 – persona\_gemini – Score 13.00 – t283\_gemini.txt}  \label{ex:f1_pos_044}
For 50 \textbf{years}, the oil industry has boomed in Scotland, but it is now changing.
This is leading to rising unemployment among workers, increased use of food banks, and losses for communities.
Ravishaan Rahel Muthiah, a digital campaigner for Greenpeace UK, is \textbf{working} for a just transition to clean energy and green \textbf{jobs}.
To understand the situation better, Greenpeace \textbf{interviewed} current and former oil workers.
Wullie spoke about the hardships of \textbf{working} offshore, as well as redundancies and the health impacts he has faced.
Robbie described the barriers to transitioning into new \textbf{work}, like the high fees for training.
Yvonne \textbf{shared} her view that the government needs to prioritize the planet over the profits of oil companies.
Alan, another worker, spoke of a hopeful future in renewables.
These \textbf{stories} highlight challenges that are common among workers. \textbf{People} tied to polluting industries should not be the \textbf{ones} to \textbf{bear} the burden of the climate crisis.
That is why Greenpeace UK, along with \textbf{Friends} of the Earth Scotland and \textbf{Platform}, is \textbf{campaigning} for supported transitions to secure green \textbf{jobs}.
To hear these community voices for yourself, \textbf{watch} the \textbf{film}.

% matched lemmas: bear, campaign, film, friend, interview, job, one, people, platform, share, story, watch, work, year
\end{textsample}
